\item \textbf{ACTIVIDAD} Si se lanzan dos dados de seis caras, los posibles resultados de sumar el número de puntos en las caras superiores oscilan entre 2 y 12. Las probabilidades de estos resultados varían según la cantidad de formas en que se puede lograr un resultado en particular. Por ejemplo, solo hay una forma de lograr un 2: 1-1. Pero hay seis formas de lograr un resultado de 7: 1-6, 2-5, 3-4, 4-3, 5-2 y 6-1. Por tanto, un 7 es seis veces más probable que un 2. El gráfico de barras en la figura TP9.2 muestra la probabilidad teórica de todos los resultados posibles para dos dados. Experimentalmente, si los dos dados se lanzaran 100 veces y se dibujara un histograma de posibles resultados, tendría una forma muy similar a este gráfico de probabilidad. 

(a) ¿Qué pasa si arrojamos \textit{tres} dados? ¿Cómo se verá el histograma? Haga esto en su grupo. Lance tres dados 100 veces y haga un histograma de los resultados. ¿Cómo se diferencia la forma del histograma resultante de la curva de probabilidad que se muestra en la figura TP9.2 para dos dados? Haga un gráfico de probabilidad teórico para tres dados como en la figura TP9.2 y compárelo con su histograma. 
(b) ¿Qué pasa si lanza los dados la cantidad de Avogadro? (¡No intente esto!) ¿Cómo se vería el histograma? (Realice una predicción de cómo varía el gráfico para tres dados respecto al de dos dados.) 
(c) Suponga que configura laboriosamente un número de Avogadro de dados sobre una mesa, con todos los dados teniendo un 1 en su cara superior. Luego sacuda la mesa por unos segundos. Cuando ha sumado todos los números en las caras superiores, ¿cuál es el resultado más probable? 
(d) Ahora imagine que sacude la mesa de nuevo. ¿Qué tan probable es que todos los dados vuelvan a tener un 1 en sus caras superiores? 
(e) ¿Qué tiene que ver todo esto con la entropía?
